\section*{Chapitre \ref{ch:dv:autocallables} --- Autocallables}

\begin{solution}{exo:dv:autocallables:1}
Une obligation zéro-coupon remboursant le pair (position longue)~; un put européen knock-in de strike 100~\% et de barrière 60~\% (position courte)~; une série d'\omterm{def:dv:barriers-and-digitals:digital}{options
digitales} versant le coupon lorsque l'indice est supérieur ou égal à 70~\% à chaque date, avec mémoire (position longue)~; la clause de remboursement anticipé, une
série de digitales à 100~\% qui mettent fin à la note (position courte, car l'émetteur ne choisit rien, mais l'investisseur perd les coupons restants
et le potentiel de hausse au-delà du pair).
\end{solution}

\begin{solution}{exo:dv:autocallables:2}
L'investisseur finance le coupon surtout en vendant le put knock-in. Sous l'effet du smile, ce put est valorisé aux volatilités élevées des
strikes faibles, donc il vaut davantage et finance un coupon plus élevé~; les volatilités à la monnaie de courte échéance, plus basses dans le smile, modifient aussi les
probabilités de remboursement anticipé. Une volatilité plate sous-évalue le put et donc le coupon qu'il permet de financer~: 3.14~\% contre 6.18~\% par an.
\end{solution}

\begin{solution}{exo:dv:autocallables:3}
$(98-95.99)/2.60=0.78\%$ par trimestre, 3.1~\% par an.
\end{solution}

\begin{solution}{exo:dv:autocallables:4}
L'émetteur est long du put knock-in. Des dividendes plus élevés abaissent le forward, ce qui rend le put plus précieux~: l'émetteur y gagne, donc il est
long de dividendes (dans l'exemple du chapitre, un rendement de 1~\% fait varier la valeur de la note de $-0.20$ pour l'investisseur et de $+0.20$ pour l'émetteur). Il vend
son exposition via des contrats à terme sur dividendes ou des swaps.
\end{solution}

\begin{solution}{exo:dv:autocallables:5}
Le delta de la note pour l'investisseur croît à mesure que le sous-jacent s'approche de la barrière, car le put knock-in gagne rapidement de la valeur à cet endroit~;
l'émetteur, en position courte sur la note, couvre ce delta en sous-jacent et achète lorsque celui-ci baisse (et vend lorsqu'il monte). Au moment du knock-in, la note
devient une position longue dont le delta vaut environ un, bien inférieur au delta juste au-dessus de la barrière, et l'émetteur vend la différence
pendant la baisse.
\end{solution}

\begin{solution}{exo:dv:autocallables:6}
Juste au-dessus de la barrière, la valeur de la note passe d'environ le pair plus le coupon à la valeur après knock-in sur une distance qui se réduit comme
$\sigma\sqrt\tau$~: moins il reste de temps, plus la chute est raide et plus le delta est grand avant le knock-in. Après celui-ci, le delta vaut environ
un, quel que soit le temps restant. L'écart, la falaise, croît comme $1/\sqrt\tau$~: 97, 254 et 497 millions pour un an, trois mois et un
mois dans l'exemple du chapitre.
\end{solution}

\begin{solution}{exo:dv:autocallables:7}
Le coupon juste passe de 1.55~\% à 1.91~\% par trimestre (7.62~\% par an), et la probabilité de knock-in de 6.3~\% à 10.4~\%~: une barrière observée
chaque jour est franchie par des trajectoires qui remontent avant l'échéance, si bien que l'investisseur vend un put plus précieux et reçoit davantage en contrepartie.
\end{solution}

\begin{solution}{exo:dv:autocallables:8}
Les couvertures sont locales~: lors d'un knock-in constaté chaque jour, le delta saute et le desk doit vendre immédiatement, sur un marché en baisse, aux côtés de tous les autres
émetteurs ayant des notes au même niveau~; un gap franchissant le niveau oblige à vendre à un prix moins favorable (chapitre~15). Le \omterm{def:dv:greeks-and-the-hedging-pnl:greeks}{vega} et le skew du book
changent aussi brusquement. L'événement exige une réserve pour gap, des limites de concentration aux niveaux de barrière et un prépositionnement.
\end{solution}

\begin{solution}{pb:dv:autocallables:1}
\textbf{1.} 100.44 (erreur type 0.06)~; probabilité de knock-in 15.6~\%.
\textbf{2.} Un put down-and-in à observation quotidienne de strike 100~\% et de barrière 75~\%, et la hausse, plafonnée à 103~\%~; ils financent le
coupon qui s'accumule.
\textbf{3.} Knock-out~: pair plus 15~\% par an jusqu'à la date de knock-out. Aucun des deux~: pair plus 30~\% à l'échéance. Knock-in sans knock-out ultérieur~:
le pair multiplié par la performance finale, plafonné au pair.
\textbf{4.} La barrière peut être franchie par une trajectoire qui remonte avant l'échéance, et elle est observée chaque jour, y compris les jours de fortes baisses
et de gaps.
\textbf{5.} Le \omterm{def:dv:autocallables:snowball}{snowball}, 15~\% contre 6.18~\% par an~: sa barrière de knock-in est plus élevée (75~\% contre 60~\%) et observée quotidiennement, si bien que l'investisseur vend
un put bien plus précieux.
\textbf{6.} 7.63.
\textbf{7.} 1.00.
\textbf{8.} $(7.63-1.00)\times0.75\times100$ millions $=497$ millions.
\textbf{9.} 254 millions et 97 millions.
\textbf{10.} Achats du sous-jacent, de plus en plus importants à mesure que le delta croît, passant d'environ un à 100~\% à 7.6 près de la barrière.
\textbf{11.} Une barrière dix fois plus grande~: une vente de l'ordre de cinq milliards concentrée sur un seul niveau.
\textbf{12.} Environ 0.6 fois le notionnel avec douze mois restants et environ 6 fois avec un mois.
\textbf{13.} Ce sont les couvertures les moins chères et les plus rapides pour un indice, et c'est là que le flux est exécuté en premier~; la base et l'activité le montrent.
\textbf{14.} Limiter l'émission sur un même sous-jacent en fonction de son volume de transactions et du delta maximal que la couverture peut atteindre.
\textbf{15.} Notionnel en circulation par sous-jacent, distribution des niveaux de knock-in et des échéances, et gamma des dealers des émetteurs.
\textbf{16.} Une réserve pour le gap et pour l'impact de marché de la vente~: dimensionner la vente, estimer son coût à partir de la profondeur et des modèles d'impact
(livre~8), et la provisionner, pondérée par la probabilité d'atteindre la barrière à l'approche de l'échéance.
\textbf{17.} Réduire le delta en amont de la barrière (sous-couvrir les derniers points), acheter des put spreads autour de la barrière, ou prévendre une partie du
delta de knock-in de façon contrôlée.
\textbf{18.} L'investisseur, dont la note perd désormais avec l'indice~; le P\&L de l'émetteur dépend de la qualité d'exécution de la vente.
\textbf{19.} 497 millions de sous-jacent (au niveau de la barrière) pour 100 millions de notes~: le delta passe de 7.63 à 1.00.
\textbf{20.} De nombreux émetteurs couvrent des notes ayant la même barrière et les mêmes échéances, si bien que la vente de couverture à ce niveau est celle de l'ensemble du marché.
\end{solution}

\begin{solution}{iq:dv:autocallables:1}
Une obligation plus~: vente d'un put knock-in (le capital à risque), achat de digitales de coupon (avec mémoire dans un Phoenix), vente de la hausse via le
call anticipé. Le coupon est financé par la prime du put et par le call qui écourte la note dans les scénarios favorables.
\iqlookfor{la décomposition et qui est vendeur de quoi.}
\end{solution}

\begin{solution}{iq:dv:autocallables:2}
Commencer par la \omterm{def:dv:local-volatility:model}{volatilité locale}, qui s'ajuste aux options vanille qui déterminent les distributions terminales~; vérifier avec un modèle local-stochastique pour les
coupons et calls dépendant du skew forward~; constituer une réserve pour l'écart. La volatilité plate sous-évalue le put knock-in et valorise mal les
digitales~: elle divise à peu près par deux le coupon juste dans l'exemple du chapitre.
\iqlookfor{la \omterm{def:dv:local-volatility:model}{volatilité locale}, ses limites et le sens de l'erreur liée à la volatilité plate.}
\end{solution}

\begin{solution}{iq:dv:autocallables:3}
Position longue en \omterm{def:dv:greeks-and-the-hedging-pnl:greeks}{vega} (court terme et longue échéance), position longue en skew, position longue en dividendes et, pour les worst-of, position courte en corrélation. L'émetteur vend de la volatilité (options,
\omterm{def:dv:variance-swaps-and-volatility-derivatives:vs}{swaps de variance}), du skew (risk reversals) et des dividendes (contrats à terme, swaps) au marché, et conserve la corrélation ou la compense avec des indices face à
des options sur actions individuelles.
\iqlookfor{les quatre expositions et les instruments de recyclage.}
\end{solution}

\begin{solution}{iq:dv:autocallables:4}
Utiliser des nombres aléatoires communs pour les prix après choc~; de nombreuses trajectoires~; lisser les caractéristiques digitales (remplacer la digitale par un call spread étroit, comme le ferait une
surcouverture) ou utiliser des estimateurs pathwise et par rapport de vraisemblance (chapitre~23)~; regrouper les \omterm{def:dv:greeks-and-the-hedging-pnl:greeks}{grecques} par compartiments plutôt que de choquer chaque nœud~; vérifier la stabilité
selon les graines.
\iqlookfor{les nombres aléatoires communs et le traitement des discontinuités.}
\end{solution}

\begin{solution}{iq:dv:autocallables:5}
Agréger le notionnel en circulation par sous-jacent, niveau de knock-in et échéance~; pour chaque niveau, calculer la variation du delta des émetteurs lorsqu'il est
franchi (la falaise) en fonction du temps restant jusqu'à l'échéance~; comparer avec le volume quotidien et la profondeur du sous-jacent et des contrats à terme.
\iqlookfor{notionnel multiplié par la concentration par barrière et par échéance, multiplié par le gamma des dealers.}
\end{solution}

\begin{solution}{iq:dv:autocallables:6}
Estimer la vente du marché à ce niveau~; réduire mon propre delta en amont et acheter une protection (puts ou put spreads sous le niveau)~; élargir les cotations
pour les nouvelles notes~; se préparer à fournir de la liquidité après la rupture du niveau si la vente déborde~; prévenir la direction des risques et dimensionner la réserve pour gap.
\iqlookfor{anticiper le flux, la protection et la liquidité après la rupture.}
\end{solution}
